\section{Learning from Experience}

\subsection{Experience-Driven Self-Improvement: Evaluation Design}
\label{self_improvement_evaluation}

Beyond process quality within a single run, practical automated research requires agents to improve as they accumulate experience over extended workflows.
We evaluate this evolving capability at two scales: \textbf{intra-task self-improvement} tests whether experience from earlier iterations improves later solutions to the same task, while \textbf{inter-task self-improvement} measures whether experience from solved tasks improves performance on a held-out task.

$\mathbf{M_{\mathrm{intra}}}$\textbf{: Intra-Task Self-Improvement.} Intra-task self-improvement evaluates whether an agent can leverage experience from earlier iterations of the same task to propose better solutions later on.
This is crucial for automated research, where solving a task typically requires iterative exploration rather than a single common-sense guess.

To isolate the effect of this accumulated experience, we adopt a counterfactual design that compares the quality of a single solution the model proposes with and without experience.
From the agent's trajectory, we select a \textbf{branch point} from which two conditions continue optimizing the same intermediate solution.
In the with-experience condition, the agent continues normally with its accumulated experience retained.
In the without-experience condition, we re-initialize the agent and erase its prior context, on-disk notes, and in-code comments while preserving the solution at the branch point.
We compare the next commit produced under the two conditions, and their gap measures the degree to which the proposal relies on intra-task experience.
Notably, we focus on the first commit after the branch point because further iteration may reconstruct the erased experience, obscuring its isolated effect.
Let $S^{\text{exp}}$ and $S^{\text{no\_exp}}$ be the scores of the first commit after the branch point under the with- and without-experience conditions.
The gain is their difference,
\[
    \Delta S_{\text{intra}} = S^{\text{exp}} - S^{\text{no\_exp}} \in [-1,+1].
\]
A larger gap means the quality of the model's proposed solution depends more heavily on its prior exploration experience, whereas a smaller gap means it depends less on that experience.

$\mathbf{M_{inter}}$\textbf{: Inter-Task Self-Improvement.}
Complementing the intra-task setting, inter-task self-improvement evaluates whether an agent can extract reusable experience from a solved source task and apply it to a held-out target task, capturing its capacity for continued improvement across sustained auto research workflows.

Specifically, given a model and a source--target task pair, the model extracts lessons from its completed source trajectory and then attempts the target under two conditions: a baseline run without lessons and an augmented run with them.
We hold the model and all target-task conditions fixed, including the harness, execution environment, and resource limits, and use separate workspaces so that the augmented run receives only the extracted lessons, not source-task artifacts.
If their scores are $S^{(0)}$ and $S^{(+)}$, respectively, the transfer gain
\[
    \Delta S_{\text{inter}} = S^{(+)} - S^{(0)} \in [-1,+1]
\]
provides a direct measure of whether the model can improve target performance by extracting transferable experience and applying it effectively.


\subsection{Experience-Driven Self-Improvement: Results}

\subsubsection{Intra-Task Experience Reuse}
\label{sec:m1_in_task}



We evaluate intra-task self-improvement by measuring how much the experience an agent accumulates within a single run improves the next solution it produces, following the counterfactual design of \S\ref{self_improvement_evaluation}.
For trajectories lacking an available commit after the branch point, we drop the corresponding task for all models to ensure a fair comparison, retaining 32 tasks in total.


 
\textbf{Experience erasure.}
For each retained trajectory, we place the branch point near the midpoint of the run, late enough for the agent to have accumulated meaningful experience yet early enough to leave headroom for that experience to make a measurable difference.
To erase the experience, we re-initialize Claude Code from scratch, clearing both its in-context history and any notes it persisted to disk.
Since some models leave prior findings as code comments, we additionally strip all comments.
Finally, only the solution at the branch point is carried over, so the two conditions optimize \textbf{the same} starting solution.

\begin{wrapfigure}{r}{0.52\linewidth}
\vspace{-\baselineskip}
\centering
\includegraphics[width=\linewidth]{Figures/m1_in_task_summary.pdf}
\caption{Per-model first-commit score with and without retained experience (bars, left axis) and the corresponding intra-task gain $\Delta$ (line, right axis), averaged over $32$ retained trajectories.}
\label{fig:m1_in_task}
\vspace{-\baselineskip}
\end{wrapfigure}


Figure~\ref{fig:m1_in_task} reports both the first-commit scores after the branch point and the corresponding intra-task gain for each evaluated model.

\textbf{Intra-task experience generally improves the next commit across models.}
The sole exception, Kimi-K2.7-Code ($-0.0127$), is driven by a small number of retained-experience trajectories in which incomplete intermediate proposals receive zero scores, lowering the overall mean gain. 
Nevertheless, Kimi still benefits from experience on more tasks than it is harmed ($17$ vs.\ $10$; Appendix Figure~\ref{fig:m1_in_task_outcomes}).
The complete task-level sign counts show the same tendency for all seven models across the $32$ tasks.

\textbf{Models differ widely in how much they rely on intra-task experience.}
Opus-4.7 records the smallest positive gain ($+0.0362$), possibly because its top Solution Framing (C1) score in \S\ref{sec:process_results} allows it to formulate strong solutions with little support from prior exploration.
By contrast, several models with lower overall performance, including DeepSeek-V4-Pro, Gemini-3.1-Pro, and LongCat-2.0, show substantially larger gains, suggesting that accumulated experience has a stronger influence on their next commits.
LongCat provides the clearest example, combining the largest gain ($+0.1454$) with one of the weakest solution framing, such that much of its next-commit quality depends on the experience accumulated before the branch point.
These results generally suggest that weaker models tend to rely more heavily on experience accumulated through multi-step exploration to improve solution quality.


\textbf{Why retained experience is usually beneficial, and when it backfires.}
Our trajectory case studies (Appendix~\S\ref{sec:in_task_traj}) include both positive and negative examples, helping clarify how retained experience shapes the next commit.
On the positive side, accumulated experience enables the next commit to avoid known dead ends, reuse tuned configurations, and carry forward hard-won implementations, thereby improving the commit quality.
However, in a smaller number of cases, retained experience backfires: the carried-over state may preserve a premature or misleading conclusion, or anchor the agent to a local optimum.
These findings suggest that current models still leave room to improve in how reliably they exploit their own experience.


\subsubsection{Inter-Task Experience Reuse}
\label{sec:m1_inter_task}

We reuse the three lesson-free rollouts from the outcome-level evaluation (\S\ref{sec:main_results}) as each model's baseline, yielding scores $S^{(0)}$.
Based on baseline trajectory quality, we select one source task from each AutoLab category, requiring both strong outcomes and substantive exploration; the resulting four source tasks are shared across all evaluated models.
For each source, every model extracts lessons from its own best baseline trajectory and records them in a concise \texttt{lessons.md} file summarizing what worked, what failed, and general recommendations that may transfer to unseen tasks; Appendix~\ref{app:inter_task_lessons_example} provides a representative example.
Among the remaining 32 tasks, we retain 19 whose baseline performance leaves every model sufficient room to improve, pair each with the source from its category, and fix the resulting source--target pairs across all models.
Each model then performs three new rollouts on each target in isolated workspaces, receiving only its own lessons from the paired source and yielding scores $S^{(+)}$.
Appendix~\ref{app:m1_inter_task_protocol} provides the exact selection rules and task lists.

Figure~\ref{fig:m1_inter_task} compares avg@3 with and without trajectory-derived experience and reports the corresponding inter-task gains; best@3 exhibits a similar overall pattern and is reported in Appendix~\ref{app:m1_inter_task_best}.


\textbf{Initial performance does not reliably predict a model's ability to improve through experience.}
Several leading models nevertheless show clear strengths: GPT-5.5 and GLM-5.2 improve under both metrics, with GPT gaining more on avg@3 than best@3 ($+0.063$ vs.\ $+0.022$), reflecting broader gains across runs, and GLM showing the reverse ($+0.040$ vs.\ $+0.067$), driven by larger improvements in its best runs.
Opus-4.7 is nearly unchanged on avg@3 ($+0.001$) but improves on best@3 ($+0.038$), showing that experience can raise its best-achieved performance without changing its average.
However, strong initial performance is neither necessary nor sufficient for effective reuse: DeepSeek-V4-Pro has the weakest lesson-free baseline yet records the largest gains ($+0.093$ on avg@3 and $+0.071$ on best@3), whereas the higher-performing Gemini-3.1-Pro declines on avg@3 ($-0.017$) and remains unchanged on best@3 ($+0.003$).
This distinction matters in sustained auto research workflows: as agents accumulate experience across tasks, performance gaps may narrow or widen, and initially lower-performing models may eventually overtake those that start ahead.

\begin{wrapfigure}{r}{0.52\linewidth}
\vspace{-\baselineskip}
\centering
\includegraphics[width=\linewidth]{Figures/m1_inter_task_reward_gap.pdf}
\caption{Per-model avg@3 with and without trajectory-derived experience (bars, left axis) and the corresponding inter-task gain (line, right axis).}
\label{fig:m1_inter_task}
\vspace{-\baselineskip}
\end{wrapfigure}

\textbf{Experience reuse can improve performance but remains unstable: successful transfer abstracts general principles, whereas failures misapply source-specific tactics or reinforce evaluator-specific shortcuts.}
Although most models show a positive aggregate gain, transfer remains mixed at the task level, improving performance on some targets while reducing it on others (Appendix Figure~\ref{fig:m1_inter_task_outcomes}).
DeepSeek-V4-Pro's lessons emphasize constraint checking, verification, and rollback, directly addressing Feedback Control (C3), its weakest dimension in our process evaluation (\S\ref{sec:process_results}).
Its zero-score outcomes fall from 13 of 57 lesson-free rollouts to none with lessons, helping explain its large aggregate gain.
By contrast, Opus-4.7 spends six rounds applying a source-derived caching tactic to mostly unique Levenshtein inputs, where caching adds overhead rather than reducing computation.
Gemini-3.1-Pro reveals a deeper risk: after extracting ``semantic mocking'' as transferable knowledge, it caches a SHA-256 digest during warmup and returns it during timed evaluation, producing an apparent $+0.620$ best@3 gain without accelerating SHA-256 itself.

\textbf{Experience transfers more effectively through explicitly extracted, self-generated lessons.}
To further investigate effective strategies for experience reuse, we vary the main design along two axes: representation, comparing extracted lessons with access to the full source workspace, and source, comparing self-generated lessons with those produced by another model.
For representation, explicitly extracted lessons outperform access to raw source workspaces for all three tested models under both metrics, suggesting that lesson extraction improves transfer by filtering noise and surfacing transferable knowledge.
For source, self-generated lessons outperform cross-model lessons for both GLM-5.2 and LongCat-2.0: lessons that improve the stronger GLM do not benefit LongCat, while GLM loses its self-reuse gains when using LongCat's lessons, showing that lesson effectiveness depends on compatibility with the receiving model rather than producer strength alone.
Complete experimental setup and results are reported in Appendix~\ref{sec:reuse_forms}.

Overall, our results show that even a single transfer step can improve subsequent task performance, highlighting the potential of experience reuse for cumulative improvement over longer auto research workflows.
However, the unstable gains and varying effectiveness of reuse strategies indicate that reliable long-horizon self-improvement requires better mechanisms throughout the experience-reuse pipeline, from extracting and selecting transferable lessons to adapting, applying, and revising them in response to feedback.
